% Part: lambda-calculus
% Chapter: introduction
% Section: unique-readability

\documentclass[../../../include/open-logic-section]{subfiles}

\begin{document}

\olfileid[pt]{lam}{syn}{unq}
\olsection{Legibilidade única}

Podemos nos perguntar se há, para cada termo, uma única maneira de formá-lo,
e há. Para cada termo lambda, há apenas uma maneira de construí-lo e
interpretá-lo. Na discussão seguinte, uma \emph{formação} é o procedimento
de construir um termo usando as regras de formação (uma ou várias vezes) de
\olref[trm]{defn:term}.

\begin{lem}\ollabel{lem:term-start}
Um termo começa por uma variável ou por um parêntese.
\end{lem}

\begin{proof}
Algo conta como termo somente se for construído de acordo com
\olref[trm]{defn:term}. Se for resultado de
\olref[trm]{defn:term-var}, deve ser uma variável. Se for resultado de
\olref[trm]{defn:term-abs} ou \olref[trm]{defn:term-app}, começa por um
parêntese.
\end{proof}

\begin{lem}\ollabel{lem:app-start}
O resultado de uma aplicação começa por dois parênteses ou por um parêntese
e uma variável.
\end{lem}

\begin{proof}
Se $M$ é resultado de uma aplicação, tem a forma $(PQ)$ e, portanto, começa
por um parêntese. Como $P$ é um termo, por \olref{lem:term-start} ele começa
por um parêntese ou por uma variável.
\end{proof}

\begin{lem}\ollabel{lem:initial}
Nenhuma parte inicial própria de um termo é, ela própria, um termo.
\end{lem}

\begin{prob}
Prove \olref[lam][syn][unq]{lem:initial} por indução no comprimento dos
termos.
\end{prob}

\begin{prop}[Legibilidade única] \ollabel{prop:unq}
Há uma formação única para cada termo. Em outras palavras, se um termo $M$
é formado por uma formação, então essa é a única formação capaz de formar
esse termo.
\end{prop}

\begin{proof}
  Provamos isso por indução na formação dos termos.
  \begin{enumerate}
    \item $M$ tem a forma $x$, onde $x$ é alguma variável. Como os resultados
      de abstrações e aplicações sempre começam por parênteses, elas não
      podem ter sido usadas para construir $M$; portanto, a formação de $M$
      deve ser um único passo de
      \olref[trm]{defn:term}\olref[trm]{defn:term-var}.
    \item $M$ tem a forma $(\lambd[x][N])$, onde $x$ é alguma variável e
      $N$ é um termo. Ele não poderia ter sido construído de acordo com
      \olref[trm]{defn:term}\olref[trm]{defn:term-var}, pois não é uma única
      variável. Por \olref{lem:app-start}, ele não é resultado de uma
      aplicação. Logo, $M$ só pode ser resultado de uma abstração sobre~$N$.
      Pela hipótese de indução, sabemos que a formação de $N$ é única.
    \item $M$ tem a forma $(PQ)$, onde $P$ e $Q$ são termos. Como começa por
      um parêntese, não pode também ser construído por
      \olref[trm]{defn:term}\olref[trm]{defn:term-var}. Por
      \olref{lem:term-start}, $P$ não pode começar por $\lambd$; logo,
      $(PQ)$ não pode ser resultado de uma abstração. Suponha agora que
      houvesse outra maneira de construir $M$ por aplicação, por exemplo,
      que ele também tivesse a forma $(P'Q')$. Então $P$ seria um segmento
      inicial próprio de $P'$ (ou vice-versa), o que é impossível por
      \olref{lem:initial}. Assim, $P$ e $Q$ são determinados de modo único,
      e pela hipótese de indução sabemos que as formações de $P$ e $Q$ são
      únicas.
  \end{enumerate}
\end{proof}

Uma paráfrase mais legível da proposição acima é a seguinte:
\begin{prop}
  Um termo $M$ só pode ter uma das formas seguintes:
  \begin{enumerate}
    \item $x$, onde $x$ é uma variável determinada de modo único por $M$.
    \item $(\lambd[x][N])$, onde $x$ é uma variável e $N$ é outro termo,
      ambos determinados de modo único por $M$.
    \item $(PQ)$, onde $P$ e $Q$ são dois termos determinados de modo único
      por $M$.
  \end{enumerate}
\end{prop}

\end{document}

